\documentclass[11pt]{article}

\usepackage[margin=1in]{geometry}
\usepackage{amsmath,amssymb,amsthm,mathtools}
\usepackage{enumitem}
\usepackage{booktabs}
\usepackage{tabularx}
\usepackage{microtype}
\usepackage[hidelinks]{hyperref}
\usepackage{xurl}

\newtheorem{definition}{Definition}
\newtheorem{lemma}{Lemma}
\newtheorem{remark}{Remark}

\title{Challenge Routes and Audience Profiles for Successor Maintenance:\\[0.5ex]
\large Audience-Keyed Start Objects and Escalation Order}
\author{}
\date{Unified archive note v0.05 (2026-03-21; archive v488)}

\begin{document}
\maketitle

\begin{abstract}
A downstream citation map answers a narrow question: what is the smallest sufficient object to cite?
Later papers often need a different answer: where should a release note, a referee, or an auditor actually start, and how should that route escalate if challenged?
This note isolates that audience-facing routing layer.
It defines an \emph{audience profile} and an \emph{audience-keyed challenge route} whose start object may equal the minimal citation target or may begin one layer left for stricter review, while the companion stop-profile note separately owns the approved terminal fields.
The result is a compact, auditable home for start-here / escalate-here guidance without overloading either the citation map, the stop semantics, or the wording note.
\end{abstract}

\paragraph{Non-redundancy note.}
Synthesis~31 owns the downstream normal form itself, Synthesis~34 owns the summary-stage status-envelope / lineage-notice split, Synthesis~35 owns pair-anchored sentence reuse, Synthesis~37 owns question-keyed stop profiles, Synthesis~38 owns piece-keyed challenge branches, and Synthesis~17 owns the concrete worked bundle. This note owns only audience-keyed entry paths and coarse escalation order.

\paragraph{Scope.}
This is not a new maintenance stage, new verdict, or new receipt object. It is a publication-interface note about which already-owned object a given reader should inspect first.

\paragraph{Imports.}
The downstream suffix is imported from Synthesis~31.\cite{series:downstreamnormalforms}
Summary-stage wrapper semantics are imported from Synthesis~34.\cite{series:statusenvelopes}
Sentence-reuse policy is imported from Synthesis~35.\cite{series:sentencelocks}
The concrete worked instantiation is imported from Synthesis~17.\cite{series:workedexample}

\paragraph{Exports.}
This note exports (i) audience profiles, (ii) audience-keyed challenge routes, and (iii) a minimal-start rule tying those routes back to the citation map and the companion stop profiles.

\section{Why the routing layer should be separate}
The minimal citation target is not always the best starting point for every reader.
A release note may start from the outward-facing lineage notice.
A referee may prefer to begin one layer left at the support map or the row-level closure ledger.
An auditor may begin at the narrow status token or the numeric compare object that already has replay-facing meaning.
Those are routing differences, not wording differences and not new evidence.
So they deserve one small home of their own.

\begin{definition}[Audience profile]
An \emph{audience profile} is a pair
\[
A=(a,R_a)
\]
consisting of one reader class label $a$ and a finite default route-key set $R_a$ naming the question families that audience most often asks first.
Typical audience labels are \textsf{release\_note}, \textsf{referee}, and \textsf{auditor}.
\end{definition}

\begin{definition}[Question-keyed challenge route]
For one concrete base$\to$successor comparison, an \emph{audience-keyed challenge route}
\[
\Gamma(a,q)=(m,s,e,p)
\]
consists of one audience label $a$, one downstream question key $q$, one minimal citation target $m$, one declared start object $s$, one finite escalation list $e$, and one companion stop-profile key $p$.
The route is read operationally: start at $s$, escalate through $e$ only as needed, and consult the companion stop profile keyed by $p$ when the challenge becomes specific enough to ask where escalation may terminate.
\end{definition}

\begin{definition}[Minimal-start rule]
A challenge route satisfies the \emph{minimal-start rule} if its minimal citation target $m$ agrees with the companion citation map for the same question key.
Its declared start object $s$ may equal $m$ or may begin at a strictly narrower object when the audience is expected to want stronger review by default.
\end{definition}

\begin{definition}[Route import rule]
A challenge route satisfies the \emph{route import rule} if its stop-profile key $p$ names a companion stop profile whose question family agrees with $q$ and whose approved terminal fields are not redefined inline in the route itself.
\end{definition}

\begin{lemma}[Audience-aware routing is audit-safe]
Assume every route for a concrete successor comparison satisfies the minimal-start rule and the route import rule.
Then different audiences may begin at different objects without changing the underlying evidence or inventing a new owner for any token, field, or sentence.
\end{lemma}

\begin{proof}
The minimal-start rule ties every route back to the same citation map question family, so the routing layer cannot silently replace the smallest sufficient object.
The declared start object merely allows a stricter audience to begin one layer left when that is operationally helpful.
The route import rule keeps approved terminal fields in the companion stop-profile object rather than letting every audience route redefine them ad hoc.
So different routes change navigation, not semantics.
\end{proof}

\begin{table}[t]
\centering
\small
\begin{tabularx}{\linewidth}{@{}l l X@{}}
\toprule
Audience & Typical start & Why this is different from citation minimization \\
\midrule
Release note & outward-facing wrapper & Starts where public prose already exists. \\
Referee & support map / row ledger & Often begins one layer left so justification is visible immediately. \\
Auditor & narrow token / compare object & Starts where replay or row-checking already has operational meaning. \\
\bottomrule
\end{tabularx}
\caption{Routing is about where a reader should begin, not about which object is the smallest sufficient citation target in the abstract.}
\label{tab:audiences}
\end{table}

\section{Worked example pointer}
The maintained worked bundle now ships \path{example_successor_challenge_routes.json}.\cite{series:workedexample}
It records six routes for the canned Case-B successor: release-note entry paths for public status and compact diff, referee entry paths for public status and package completeness, and auditor entry paths for public-status token and compact diff.
Some routes start at the minimal citation target; others intentionally begin one layer left, e.g., the referee package-completeness route starts from the closure ledger rather than the compressed closure token.
Each route now points at the companion \path{example_successor_challenge_stop_profiles.json} object rather than redefining approved stop fields inline. The companion citation map still owns which minimal object corresponds to each question key, the companion stop-profile note owns which terminal fields are allowed to answer the contested semantic pieces of that question family, and the new companion \path{example_successor_challenge_branches.json} object owns the exact piece-keyed escalation walk once a specific semantic piece is contested.\cite{series:downstreamnormalforms,series:challengestops,series:challengebranches}

\paragraph{Why this helps the archive.}
The series can now keep the citation map small, the normal-form note about wording, and the sentence-lock note about reuse boundaries.
Audience-facing navigation no longer has to be improvised in release wrappers or referee notes, and the route note no longer has to pretend its coarse escalate-if-challenged list is already the exact piece-specific proof trace.


\paragraph{A minimal public challenge-route card.}
\begin{table}[t]
\centering
\small
\begin{tabularx}{\linewidth}{@{}lX@{}}
\toprule
Field & Minimal public answer \\
\midrule
Release-note public-status route & \nolinkurl{release_note_public_status_sentence} starts and minimally cites \nolinkurl{example_successor_lineage_notice.json}, escalates through \nolinkurl{example_successor_claim_support_map.json}, \nolinkurl{example_successor_continuity_verdict.json}, and \nolinkurl{example_publication_closure_verdict.json}, and imports the stop profile \nolinkurl{public_status_sentence} \\
Referee package-completeness route & \nolinkurl{referee_package_completeness} starts at \nolinkurl{example_release_closure_ledger.json}, minimally cites \nolinkurl{example_publication_closure_verdict.json}, escalates through \nolinkurl{example_publication_closure_verdict.json} and \nolinkurl{example_release_obligation_profile.json}, and imports the stop profile \nolinkurl{package_completeness} \\
Auditor compact-diff route & \nolinkurl{auditor_compact_diff} starts and minimally cites \nolinkurl{example_compare_report.json}, escalates through \nolinkurl{example_verifier_report.json} and \nolinkurl{example_successor_continuity_verdict.json}, and imports the stop profile \nolinkurl{compact_diff_summary} \\
Selection rule & Stop here only while the live issue is which audience/question pair still forces the same declared start object and coarse escalation order; do not widen to semantic-piece stop semantics or exact contested-piece walks unless those narrower owners have become the honest moving part \\
Left/right boundary & Reopen Synthesis~31 or Synthesis~35 when the live issue is the minimal citation target or sentence-reuse policy itself; widen to Synthesis~37 when the contested semantic piece or approved terminal field is live, and to Synthesis~38 once the exact piece-keyed escalation walk is the honest moving part \\
Paired-terminal boundary & Once the maintained audience/question route is fixed, later terse papers should usually stop at Synthesis~54 together with Synthesis~74--75 rather than reopen this note merely to restate a stable answer-side handoff or paired cutover judgment \\
\bottomrule
\end{tabularx}
\caption{A minimal public challenge-route card. This is the smallest routing object later terse papers can quote directly when the live issue is where one audience/question pair should begin and what coarse escalation order it inherits before any semantic-piece split is live.}
\label{tab:minimal-public-challenge-route-card}
\end{table}

This card is intentionally non-decorative: it pins the actual worked audience/question route keys, start objects, minimal targets, and imported stop-profile bindings without silently re-owning the narrower semantic-piece stop semantics or exact piece-keyed branch walks that neighboring notes already own.\cite{series:workedexample}

\paragraph{A forced public challenge-route matrix.}
\begin{table}[t]
\centering
\small
\begin{tabularx}{\linewidth}{@{}lX@{}}
\toprule
Field & Forced public answer \\
\midrule
Release-note public-status floor & \nolinkurl{release_note_public_status_sentence} still forces the same route only while the same audience/question pair agrees, the same start path and minimal path stay at \nolinkurl{example_successor_lineage_notice.json}, the same three escalation objects remain listed in the same coarse order, and the imported stop-profile key remains \nolinkurl{public_status_sentence} \\
Referee package-completeness floor & \nolinkurl{referee_package_completeness} still forces the same route only while the same audience/question pair agrees, the same start path remains \nolinkurl{example_release_closure_ledger.json}, the minimal path remains \nolinkurl{example_publication_closure_verdict.json}, the same escalation pair remains listed, and the imported stop-profile key remains \nolinkurl{package_completeness} \\
Auditor compact-diff floor & \nolinkurl{auditor_compact_diff} still forces the same route only while the same audience/question pair agrees, the same start path and minimal path stay at \nolinkurl{example_compare_report.json}, the same escalation pair remains \nolinkurl{example_verifier_report.json} then \nolinkurl{example_successor_continuity_verdict.json}, and the imported stop-profile key remains \nolinkurl{compact_diff_summary} \\
Imported-owner floor & Those shortcuts remain honest only while the route still imports the same stop-profile owner for semantic-piece stop semantics and the same branch owner for exact contested-piece walks; the route does not get to keep the same answer once either neighboring owner changes under it \\
Staleness trigger floor & This matrix goes stale as soon as the audience label changes, the question family changes, the declared start object changes, the minimal citation target changes, an escalation object is added / removed / reordered, or the route no longer binds the same stop-profile key for that audience/question pair \\
Escalation pointer & Stop here only while one maintained audience/question pair still forces the same start object and coarse escalation order; widen to Synthesis~37 once the live issue is which semantic piece or approved terminal field is enough, and widen to Synthesis~38 once the exact contested-piece walk rather than the coarse route is the honest moving part \\
\bottomrule
\end{tabularx}
\caption{A forced public challenge-route matrix. This is the smallest routing object later terse papers can quote directly when the live issue is whether one maintained audience/question pair still forces the same start object and coarse escalation order rather than merely which route exists in the abstract.}
\label{tab:forced-public-challenge-route-matrix}
\end{table}

This matrix is intentionally non-decorative: it pins the actual worked route keys, start objects, coarse escalation objects, and imported stop-profile bindings that still force one maintained route, together with the conditions under which that forcing remains honest or goes stale, without silently re-owning the narrower stop-profile or exact branch-map objects.\cite{series:workedexample}

\paragraph{One tiny challenge-route drill.}
For the worked release-note public-status question, \nolinkurl{release_note_public_status_sentence} still forces the same honest start at \nolinkurl{example_successor_lineage_notice.json}. If the challenge becomes ``which half of that wrapper claim is live?'', the route is no longer the honest stopping point and a later terse paper must widen to Synthesis~37 for the semantic-piece split and approved stop fields. If the live issue instead stays at audience-level routing but the maintained start object moves left to \nolinkurl{example_successor_claim_support_map.json} or the escalation list changes order, then the route matrix has gone stale and this note must be reopened rather than quoted as if the old routing contract still held.\cite{series:workedexample}

\begin{thebibliography}{99}\small

\bibitem{series:downstreamnormalforms}
Anonymous.
\newblock Downstream Normal Forms for Successor Maintenance: Summary, Support, Citation, Quote, and Clause Discipline.
\newblock Series synthesis note (Synthesis~31), 2026. (This archive.)

\bibitem{series:statusenvelopes}
Anonymous.
\newblock Successor Status Envelopes and Lineage Notices for Release Maintenance: Narrow Status Tokens, Reader Pointer Order, and Public Answers.
\newblock Series synthesis note (Synthesis~34), 2026. (This archive.)

\bibitem{series:sentencelocks}
Anonymous.
\newblock Sentence Locks and Inline Provenance for Successor Maintenance: Pair-Anchored Reuse, Recheck Guards, and Visible Capsules.
\newblock Series synthesis note (Synthesis~35), 2026. (This archive.)

\bibitem{series:challengestops}
Anonymous.
\newblock Challenge Stop Profiles for Successor Maintenance: Question Families, Semantic Pieces, and Approved Terminal Fields.
\newblock Series synthesis note (Synthesis~37), 2026. (This archive.)

\bibitem{series:challengebranches}
Anonymous.
\newblock Challenge Branch Maps for Successor Maintenance: Piece-Keyed Escalation from Audience Routes to Approved Terminal Fields.
\newblock Series synthesis note (Synthesis~38), 2026. (This archive.)

\bibitem{series:challengeanswerreviewmenus}
Anonymous.
\newblock Challenge Answer Review Menus for Successor Maintenance: Answer-Side Terminal Citation Normal Form.
\newblock Series synthesis note (Synthesis~54), 2026. (This archive.)

\bibitem{series:pairedterminalcitationdiscipline}
Anonymous.
\newblock Paired Terminal Citation Discipline for Review Spines: Stable Answer Stops and Adjacent Request Widening.
\newblock Series synthesis note (Synthesis~74), 2026. (This archive.)

\bibitem{series:pairedterminalcutoverrules}
Anonymous.
\newblock Paired Terminal Cutover Rules for Review Spines: When Later Terse Papers Stop, Widen, or Reopen.
\newblock Series synthesis note (Synthesis~75), 2026. (This archive.)

\bibitem{series:workedexample}
Anonymous.
\newblock Worked Example for Anonymous-DHT Receipts: Minimal Public Surface, Cross-Release Diff, and Successor Interlock.
\newblock Series synthesis note (Synthesis~17), 2026. (This archive.)

\end{thebibliography}
\end{document}
